\input format.tex

%\usepackage{arydshln}
\usepackage{graphicx}
\graphicspath{{cores/}}

\begin{document}

%%%
\setlength{\arrayrulewidth}{0.5pt}
\fontsize{9.3pt}{17pt}\selectfont
\color{gray2}

\vspace*{2mm}
\begin{center}
%{\bf\sanhao 肠道菌群分布}
{\bf\sanhao 分布}
\end{center}

\medskip

\noindent
{\bf\xiaosihao Tom先生：}


\bigskip

您的肠道菌群类别总共有301种，已知细菌298种；未知细菌3种。

您的肠道菌群多样性水平
偏低，表明您的菌群平衡较差，需要注意肠道健康
（儿童和老年人由于年龄因素，肠道菌群多样性水平比一般成年人差）。

您的肠道菌群多样性、有益菌、细菌敏感性等指标存在异常
，可能增加菌群失衡的风险，不利于肠道健康、抵抗疾病、抗氧化等，需引起注意，建议持续监测肠道健康。

{
\lantxh
\fontsize{9.3pt}{9.8pt}\selectfont
\bigskip
%\begin{tctabularx}{tabularx={m{3.75cm}m{4.9cm}<{\centering}m{1.22cm}<{\centering}m{1.8cm}<{\centering}m{2cm}<{\centering}}}
\begin{tctabularx}{tabularx={m{4.8cm}m{5.2cm}<{\centering}m{4.55cm}<{\centering}}}
%\cellcolor{kelly} &
%\cellcolor{kelly} &
%\cellcolor{kelly} &
%\cellcolor{kelly} &
%\cellcolor{kelly}
&&
\\[-6pt]
\cellcolor{kelly} \raisebox{2.614pt}{\color{black70}\bf 测试项目} &
\cellcolor{kelly} \raisebox{2.614pt}{\color{black70}\bf 分布} &
%\cellcolor{kelly} \raisebox{2.614pt}{\color{black70}\bf 检测值} &
%\cellcolor{kelly} \raisebox{2.614pt}{\color{black70}\bf 参考范围} &
\cellcolor{kelly} \raisebox{2.614pt}{\color{black70}\bf 检测结果}
%{\color{black70}\bf 测试项目} & {\color{black70}\bf 分布} & {\color{black70}\bf 检测值} &{\color{black70}\bf 参考范围} & {\color{black70}\bf 检测结果}
\end{tctabularx}

\vspace*{-4.25mm}
\fontsize{8.8pt}{11pt}\selectfont
%\begin{longtable}{m{3.75cm}m{5cm}<{\centering}m{0cm}@{}m{1.22cm}<{\centering}m{1.8cm}<{\centering}m{2cm}<{\centering}}
\begin{longtable}{m{4.8cm}m{5.2cm}<{\centering}m{0cm}@{}m{4.61cm}<{\centering}}
\hline
%\parbox[c]{\hsize}{\vskip7pt 肠道菌群多样性\\用于衡量肠道菌群物种丰富度及稳定性的指标 \vskip7pt} & \parbox[c]{\hsize}{\vskip7pt\centerline{\raisebox{-1.5ex}{\includegraphics[width=5cm,height=0.3cm]{zxjunqunduoyangxing_dis.png}}}\vskip7pt}  &\hspace*{-3.047181248cm}\raisebox{-0.45ex}{\includegraphics{triangle.pdf}}&\liuhao 3.20 & \liuhao 1.92\textasciitilde 4.41 & 表明肠道菌群组成单一，可能增加菌群失衡风险 \\
%\parbox[c]{\hsize}{\vskip7pt 肠道菌群多样性\\用于衡量肠道菌群物种丰富度及稳定性的指标 \vskip7pt} & \parbox[c]{\hsize}{\vskip7pt\centerline{\raisebox{-1.5ex}{\includegraphics[width=5cm,height=0.3cm]{zxjunqunduoyangxing_dis.png}}}\vskip7pt}  &\hspace*{-3.047181248cm}\raisebox{-0.45ex}{\includegraphics{triangle.pdf}} & \begin{minipage}{4.60cm}\begin{center}{偏低 \\ \bahao 表明肠道菌群组成单一，可能增加菌群失衡风险 }\end{center} \end{minipage} \\
\parbox[c]{\hsize}{\vskip6pt 肠道菌群多样性\\用于衡量肠道菌群物种丰富度及稳定性的指标 \vskip6pt} & \parbox[c]{\hsize}{\vskip6pt\centerline{\raisebox{-1.5ex}{\includegraphics[width=5cm,height=0.3cm]{zxjunqunduoyangxing_dis.png}}}\vskip6pt}  &\hspace*{-3.047181248cm}\raisebox{-0.45ex}{\includegraphics{triangle.pdf}} & \begin{minipage}{4.60cm}\begin{center}{偏低\\ \bahao 表明肠道菌群组成单一，可能增加菌群失衡风险 }\end{center} \end{minipage} \\
\hline
%\parbox[c]{\hsize}{\vskip7pt 有益菌\\产生有益物质，抑制病原菌的生长，帮助预防疾病 \vskip7pt} & \parbox[c]{\hsize}{\vskip7pt\centerline{\raisebox{-1.5ex}{\includegraphics[width=5cm,height=0.3cm]{zxyouyijun_dis.png}}}\vskip7pt}  &\hspace*{-1.776534248cm}\raisebox{-0.45ex}{\includegraphics{triangle.pdf}}&\liuhao 17.51 & \liuhao 0\textasciitilde 26.89 & 可能导致病原菌增多，不利于肠道健康 \\
%\parbox[c]{\hsize}{\vskip7pt 有益菌\\产生有益物质，抑制病原菌的生长，帮助预防疾病 \vskip7pt} & \parbox[c]{\hsize}{\vskip7pt\centerline{\raisebox{-1.5ex}{\includegraphics[width=5cm,height=0.3cm]{zxyouyijun_dis.png}}}\vskip7pt}  &\hspace*{-1.776534248cm}\raisebox{-0.45ex}{\includegraphics{triangle.pdf}} & \begin{minipage}{4.60cm}\begin{center}{偏低 \\ \bahao 可能导致病原菌增多，不利于肠道健康 }\end{center} \end{minipage} \\
\parbox[c]{\hsize}{\vskip6pt 有益菌\\产生有益物质，抑制病原菌的生长，帮助预防疾病 \vskip6pt} & \parbox[c]{\hsize}{\vskip6pt\centerline{\raisebox{-1.5ex}{\includegraphics[width=5cm,height=0.3cm]{zxyouyijun_dis.png}}}\vskip6pt}  &\hspace*{-1.776534248cm}\raisebox{-0.45ex}{\includegraphics{triangle.pdf}} & \begin{minipage}{4.60cm}\begin{center}{偏低\\ \bahao 可能导致病原菌增多，不利于肠道健康 }\end{center} \end{minipage} \\
\hline
%\parbox[c]{\hsize}{\vskip7pt 有害菌\\产生多种毒素，危害肠道健康 \vskip7pt} & \parbox[c]{\hsize}{\vskip7pt\centerline{\raisebox{-1.5ex}{\includegraphics[width=5cm,height=0.3cm]{zxyouhaijun_dis.png}}}\vskip7pt}  &\hspace*{-4.33716758cm}\raisebox{-0.45ex}{\includegraphics{triangle.pdf}}&\liuhao 0.04 & \liuhao 0\textasciitilde 23.24 & - \\
%\parbox[c]{\hsize}{\vskip7pt 有害菌\\产生多种毒素，危害肠道健康 \vskip7pt} & \parbox[c]{\hsize}{\vskip7pt\centerline{\raisebox{-1.5ex}{\includegraphics[width=5cm,height=0.3cm]{zxyouhaijun_dis.png}}}\vskip7pt}  &\hspace*{-4.33716758cm}\raisebox{-0.45ex}{\includegraphics{triangle.pdf}} & \begin{minipage}{4.60cm}\begin{center}{偏低 \\ \bahao - }\end{center} \end{minipage} \\
\parbox[c]{\hsize}{\vskip6pt 有害菌\\产生多种毒素，危害肠道健康 \vskip6pt} & \parbox[c]{\hsize}{\vskip6pt\centerline{\raisebox{-1.5ex}{\includegraphics[width=5cm,height=0.3cm]{zxyouhaijun_dis.png}}}\vskip6pt}  &\hspace*{-4.33716758cm}\raisebox{-0.45ex}{\includegraphics{triangle.pdf}} & \begin{minipage}{4.60cm}\begin{center}{偏低 }\end{center} \end{minipage} \\
\hline
%\parbox[c]{\hsize}{\vskip7pt 细菌敏感性\\抵抗致病菌的能力 \vskip7pt} & \parbox[c]{\hsize}{\vskip7pt\centerline{\raisebox{-1.5ex}{\includegraphics[width=5cm,height=0.3cm]{zxxijunminganxing_dis.png}}}\vskip7pt}  &\hspace*{-4.520911154cm}\raisebox{-0.45ex}{\includegraphics{triangle.pdf}}&\liuhao 0.05 & \liuhao 0\textasciitilde 31.98 & 对疾病抵抗力降低，不利于抵抗疾病 \\
%\parbox[c]{\hsize}{\vskip7pt 细菌敏感性\\抵抗致病菌的能力 \vskip7pt} & \parbox[c]{\hsize}{\vskip7pt\centerline{\raisebox{-1.5ex}{\includegraphics[width=5cm,height=0.3cm]{zxxijunminganxing_dis.png}}}\vskip7pt}  &\hspace*{-4.520911154cm}\raisebox{-0.45ex}{\includegraphics{triangle.pdf}} & \begin{minipage}{4.60cm}\begin{center}{偏低 \\ \bahao 对疾病抵抗力降低，不利于抵抗疾病 }\end{center} \end{minipage} \\
\parbox[c]{\hsize}{\vskip6pt 细菌敏感性\\抵抗致病菌的能力 \vskip6pt} & \parbox[c]{\hsize}{\vskip6pt\centerline{\raisebox{-1.5ex}{\includegraphics[width=5cm,height=0.3cm]{zxxijunminganxing_dis.png}}}\vskip6pt}  &\hspace*{-4.520911154cm}\raisebox{-0.45ex}{\includegraphics{triangle.pdf}} & \begin{minipage}{4.60cm}\begin{center}{偏低\\ \bahao 对疾病抵抗力降低，不利于抵抗疾病 }\end{center} \end{minipage} \\
\hline
%\parbox[c]{\hsize}{\vskip7pt 硫酸盐还原\\产生有害气体，刺激肠道 \vskip7pt} & \parbox[c]{\hsize}{\vskip7pt\centerline{\raisebox{-1.5ex}{\includegraphics[width=5cm,height=0.3cm]{zxliusuanyanhuanyuan_dis.png}}}\vskip7pt}  &\hspace*{-2.898737906cm}\raisebox{-0.45ex}{\includegraphics{triangle.pdf}}&\liuhao 0.47 & \liuhao 0\textasciitilde 1.53 & - \\
%\parbox[c]{\hsize}{\vskip7pt 硫酸盐还原\\产生有害气体，刺激肠道 \vskip7pt} & \parbox[c]{\hsize}{\vskip7pt\centerline{\raisebox{-1.5ex}{\includegraphics[width=5cm,height=0.3cm]{zxliusuanyanhuanyuan_dis.png}}}\vskip7pt}  &\hspace*{-2.898737906cm}\raisebox{-0.45ex}{\includegraphics{triangle.pdf}} & \begin{minipage}{4.60cm}\begin{center}{偏低 \\ \bahao - }\end{center} \end{minipage} \\
\parbox[c]{\hsize}{\vskip6pt 硫酸盐还原\\产生有害气体，刺激肠道 \vskip6pt} & \parbox[c]{\hsize}{\vskip6pt\centerline{\raisebox{-1.5ex}{\includegraphics[width=5cm,height=0.3cm]{zxliusuanyanhuanyuan_dis.png}}}\vskip6pt}  &\hspace*{-2.898737906cm}\raisebox{-0.45ex}{\includegraphics{triangle.pdf}} & \begin{minipage}{4.60cm}\begin{center}{偏低 }\end{center} \end{minipage} \\
\hline
%\parbox[c]{\hsize}{\vskip7pt 抗氧化能力\\降低炎症风险，有益肠道健康 \vskip7pt} & \parbox[c]{\hsize}{\vskip7pt\centerline{\raisebox{-1.5ex}{\includegraphics[width=5cm,height=0.3cm]{zxkangyanghuanengli_dis.png}}}\vskip7pt}  &\hspace*{-3.685436972cm}\raisebox{-0.45ex}{\includegraphics{triangle.pdf}}&\liuhao 0.01 & \liuhao 0\textasciitilde 9.33 & 不利于抗氧化和抗细胞衰老功能 \\
%\parbox[c]{\hsize}{\vskip7pt 抗氧化能力\\降低炎症风险，有益肠道健康 \vskip7pt} & \parbox[c]{\hsize}{\vskip7pt\centerline{\raisebox{-1.5ex}{\includegraphics[width=5cm,height=0.3cm]{zxkangyanghuanengli_dis.png}}}\vskip7pt}  &\hspace*{-3.685436972cm}\raisebox{-0.45ex}{\includegraphics{triangle.pdf}} & \begin{minipage}{4.60cm}\begin{center}{偏低 \\ \bahao 不利于抗氧化和抗细胞衰老功能 }\end{center} \end{minipage} \\
\parbox[c]{\hsize}{\vskip6pt 抗氧化能力\\降低炎症风险，有益肠道健康 \vskip6pt} & \parbox[c]{\hsize}{\vskip6pt\centerline{\raisebox{-1.5ex}{\includegraphics[width=5cm,height=0.3cm]{zxkangyanghuanengli_dis.png}}}\vskip6pt}  &\hspace*{-3.685436972cm}\raisebox{-0.45ex}{\includegraphics{triangle.pdf}} & \begin{minipage}{4.60cm}\begin{center}{偏低\\ \bahao 不利于抗氧化和抗细胞衰老功能 }\end{center} \end{minipage} \\
\hline
\parbox[c]{\hsize}{\vskip6pt 拟杆菌属\\常见肠道菌，多数与人体互利共生，少数菌种可能会导致菌血症等 \vskip6pt} & \parbox[c]{\hsize}{\vskip6pt\centerline{\raisebox{-1.5ex}{\includegraphics[width=5cm,height=0.3cm]{zxniganjunshu_dis.png}}}\vskip6pt}  &\hspace*{-5.426071472cm}\raisebox{-0.45ex}{\includegraphics{triangle.pdf}} & \begin{minipage}{4.60cm}\begin{center}{偏低\\ \bahao 不利于肠道菌群平衡 }\end{center} \end{minipage} \\
\hline
\parbox[c]{\hsize}{\vskip6pt 布劳特氏菌属\\发酵多种植物多糖产生乙酸盐，促进肠道健康 \vskip6pt} & \parbox[c]{\hsize}{\vskip6pt\centerline{\raisebox{-1.5ex}{\includegraphics[width=5cm,height=0.3cm]{zxbulaoteshijunshu_dis.png}}}\vskip6pt}  &\hspace*{-3.962496284cm}\raisebox{-0.45ex}{\includegraphics{triangle.pdf}} & \begin{minipage}{4.60cm}\begin{center}{偏低\\ \bahao 不利于产生有益物质及肠道健康 }\end{center} \end{minipage} \\
\hline
\parbox[c]{\hsize}{\vskip6pt 柔嫩梭菌属\\发酵纤维素产生丁酸等有益物质，抑制肠道炎症，促进肠道健康 \vskip6pt} & \parbox[c]{\hsize}{\vskip6pt\centerline{\raisebox{-1.5ex}{\includegraphics[width=5cm,height=0.3cm]{zxrounensuojunshu_dis.png}}}\vskip6pt}  &\hspace*{-0.895137992cm}\raisebox{-0.45ex}{\includegraphics{triangle.pdf}} & \begin{minipage}{4.60cm}\begin{center}{偏低\\ \bahao 不利于产生有益物质及抑制肠道炎症 }\end{center} \end{minipage} \\
\hline
\parbox[c]{\hsize}{\vskip6pt 粪球菌属\\发酵多糖，能够产生丁酸、乙酸等有益物质，调节肠道环境 \vskip6pt} & \parbox[c]{\hsize}{\vskip6pt\centerline{\raisebox{-1.5ex}{\includegraphics[width=5cm,height=0.3cm]{zxfenqiujunshu_dis.png}}}\vskip6pt}  &\hspace*{-3.396076064cm}\raisebox{-0.45ex}{\includegraphics{triangle.pdf}} & \begin{minipage}{4.60cm}\begin{center}{偏低\\ \bahao 不利于产生有益物质及调节肠道环境 }\end{center} \end{minipage} \\
\hline
\parbox[c]{\hsize}{\vskip6pt 瘤胃球菌属\\帮助降解纤维素等人体不能消化的多糖，在肠道中广泛分布 \vskip6pt} & \parbox[c]{\hsize}{\vskip6pt\centerline{\raisebox{-1.5ex}{\includegraphics[width=5cm,height=0.3cm]{zxliuweiqiujunshu_dis.png}}}\vskip6pt}  &\hspace*{-4.078039256cm}\raisebox{-0.45ex}{\includegraphics{triangle.pdf}} & \begin{minipage}{4.60cm}\begin{center}{偏低\\ \bahao 不利于降解纤维素等多糖 }\end{center} \end{minipage} \\
\hline
\parbox[c]{\hsize}{\vskip6pt 颤螺菌属\\帮助抗性淀粉和脂肪消化，保持正常体重，抑制肠道炎症 \vskip6pt} & \parbox[c]{\hsize}{\vskip6pt\centerline{\raisebox{-1.5ex}{\includegraphics[width=5cm,height=0.3cm]{zxchanluojunshu_dis.png}}}\vskip6pt}  &\hspace*{-3.13941782cm}\raisebox{-0.45ex}{\includegraphics{triangle.pdf}} & \begin{minipage}{4.60cm}\begin{center}{偏低\\ \bahao 不利于保持正常体重及抑制肠道炎症 }\end{center} \end{minipage} \\
\hline
\parbox[c]{\hsize}{\vskip6pt 副拟杆菌属\\帮助消化纤维素、抗性淀粉，保护肠道，抑制肠道炎症 \vskip6pt} & \parbox[c]{\hsize}{\vskip6pt\centerline{\raisebox{-1.5ex}{\includegraphics[width=5cm,height=0.3cm]{zxfuniganjunshu_dis.png}}}\vskip6pt}  &\hspace*{-4.742683004cm}\raisebox{-0.45ex}{\includegraphics{triangle.pdf}} & \begin{minipage}{4.60cm}\begin{center}{偏低\\ \bahao 不利于抑制肠道炎症及保护肠道 }\end{center} \end{minipage} \\
\hline
\parbox[c]{\hsize}{\vskip6pt 毛螺菌属\\发酵多种糖类产生乙酸、甲酸等物质，能保护肠黏膜，抑制肠道炎症 \vskip6pt} & \parbox[c]{\hsize}{\vskip6pt\centerline{\raisebox{-1.5ex}{\includegraphics[width=5cm,height=0.3cm]{zxmaoluojunshu_dis.png}}}\vskip6pt}  &\hspace*{-3.156247232cm}\raisebox{-0.45ex}{\includegraphics{triangle.pdf}} & \begin{minipage}{4.60cm}\begin{center}{偏低\\ \bahao 不利于产生有益物质及抑制肠道炎症 }\end{center} \end{minipage} \\
\hline
\parbox[c]{\hsize}{\vskip6pt 多尔氏菌属\\肠道的主要产气菌之一，与肠易激综合征等疾病相关 \vskip6pt} & \parbox[c]{\hsize}{\vskip6pt\centerline{\raisebox{-1.5ex}{\includegraphics[width=5cm,height=0.3cm]{zxduoershijunshu_dis.png}}}\vskip6pt}  &\hspace*{-4.085997794cm}\raisebox{-0.45ex}{\includegraphics{triangle.pdf}} & \begin{minipage}{4.60cm}\begin{center}{偏低 }\end{center} \end{minipage} \\
\hline
\parbox[c]{\hsize}{\vskip6pt 链球菌属\\多数为共生菌，少数在免疫力弱时可能引起化脓性炎症、肺炎等 \vskip6pt} & \parbox[c]{\hsize}{\vskip6pt\centerline{\raisebox{-1.5ex}{\includegraphics[width=5cm,height=0.3cm]{zxlianqiujunshu_dis.png}}}\vskip6pt}  &\hspace*{-3.71153018cm}\raisebox{-0.45ex}{\includegraphics{triangle.pdf}} & \begin{minipage}{4.60cm}\begin{center}{偏低\\ \bahao 不利于肠道菌群平衡 }\end{center} \end{minipage} \\
\hline
\parbox[c]{\hsize}{\vskip6pt 普雷沃氏菌属\\肠道共生菌，偏素食人群中含量高，少数菌种与感染、牙周疾病有关 \vskip6pt} & \parbox[c]{\hsize}{\vskip6pt\centerline{\raisebox{-1.5ex}{\includegraphics[width=5cm,height=0.3cm]{zxpushijunshu_dis.png}}}\vskip6pt}  &\hspace*{-0.701270077999999cm}\raisebox{-0.45ex}{\includegraphics{triangle.pdf}} & \begin{minipage}{4.60cm}\begin{center}{偏低\\ \bahao 不利于肠道菌群平衡 }\end{center} \end{minipage} \\
\hline
\parbox[c]{\hsize}{\vskip6pt 罗斯拜瑞氏菌属\\产生丁酸等有益物质，抑制肠道炎症，有利于肠道及人体健康 \vskip6pt} & \parbox[c]{\hsize}{\vskip6pt\centerline{\raisebox{-1.5ex}{\includegraphics[width=5cm,height=0.3cm]{zxluosibolishijunshu_dis.png}}}\vskip6pt}  &\hspace*{-4.63264094cm}\raisebox{-0.45ex}{\includegraphics{triangle.pdf}} & \begin{minipage}{4.60cm}\begin{center}{偏低\\ \bahao 不利于产生有益物质及抑制肠道炎症 }\end{center} \end{minipage} \\
\hline
\parbox[c]{\hsize}{\vskip6pt 萨特氏菌属\\可能与胃肠道感染相关 \vskip6pt} & \parbox[c]{\hsize}{\vskip6pt\centerline{\raisebox{-1.5ex}{\includegraphics[width=5cm,height=0.3cm]{zxsateshijunshu_dis.png}}}\vskip6pt}  &\hspace*{-2.765235062cm}\raisebox{-0.45ex}{\includegraphics{triangle.pdf}} & \begin{minipage}{4.60cm}\begin{center}{偏低 }\end{center} \end{minipage} \\
\hline
\parbox[c]{\hsize}{\vskip6pt 双歧杆菌属\\有益菌，降解人体不能消化的多糖，产乳酸，调节免疫及肠道环境 \vskip6pt} & \parbox[c]{\hsize}{\vskip6pt\centerline{\raisebox{-1.5ex}{\includegraphics[width=5cm,height=0.3cm]{zxshuangqiganjunshu_dis.png}}}\vskip6pt}  &\hspace*{-3.643260854cm}\raisebox{-0.45ex}{\includegraphics{triangle.pdf}} & \begin{minipage}{4.60cm}\begin{center}{偏低\\ \bahao 不利于调节免疫及肠道环境 }\end{center} \end{minipage} \\
\hline
\parbox[c]{\hsize}{\vskip6pt 阿克曼氏菌属\\降解粘蛋白、调节免疫，有利于肠黏膜完整性，保持正常体重 \vskip6pt} & \parbox[c]{\hsize}{\vskip6pt\centerline{\raisebox{-1.5ex}{\includegraphics[width=5cm,height=0.3cm]{zxaikemanshijunshu_dis.png}}}\vskip6pt}  &\hspace*{-3.809675024cm}\raisebox{-0.45ex}{\includegraphics{triangle.pdf}} & \begin{minipage}{4.60cm}\begin{center}{偏低\\ \bahao 不利于调节免疫及保持正常体重 }\end{center} \end{minipage} \\
\hline
\parbox[c]{\hsize}{\vskip6pt 梭菌属\\多数为致病菌，可能引起腹泻、肠炎等疾病 \vskip6pt} & \parbox[c]{\hsize}{\vskip6pt\centerline{\raisebox{-1.5ex}{\includegraphics[width=5cm,height=0.3cm]{zxsuojunshu_dis.png}}}\vskip6pt}  &\hspace*{-3.462549602cm}\raisebox{-0.45ex}{\includegraphics{triangle.pdf}} & \begin{minipage}{4.60cm}\begin{center}{偏低 }\end{center} \end{minipage} \\
\hline
\parbox[c]{\hsize}{\vskip6pt 嗜血杆菌属\\共生菌，某些菌种在一定条件下可引起化脓性感染、败血症等 \vskip6pt} & \parbox[c]{\hsize}{\vskip6pt\centerline{\raisebox{-1.5ex}{\includegraphics[width=5cm,height=0.3cm]{zxshixueganjunshu_dis.png}}}\vskip6pt}  &\hspace*{-3.666298334cm}\raisebox{-0.45ex}{\includegraphics{triangle.pdf}} & \begin{minipage}{4.60cm}\begin{center}{偏低 }\end{center} \end{minipage} \\
\hline
\parbox[c]{\hsize}{\vskip6pt 柯林斯氏菌属\\发酵多种糖类产生氢气、乙醇等 \vskip6pt} & \parbox[c]{\hsize}{\vskip6pt\centerline{\raisebox{-1.5ex}{\includegraphics[width=5cm,height=0.3cm]{zxkelinsishijunshu_dis.png}}}\vskip6pt}  &\hspace*{-3.870431024cm}\raisebox{-0.45ex}{\includegraphics{triangle.pdf}} & \begin{minipage}{4.60cm}\begin{center}{偏低\\ \bahao 不利于发酵糖类 }\end{center} \end{minipage} \\
\hline
\parbox[c]{\hsize}{\vskip6pt 假单胞菌属\\分布广泛，部分菌种为条件致病菌，引起医院内获得性感染 \vskip6pt} & \parbox[c]{\hsize}{\vskip6pt\centerline{\raisebox{-1.5ex}{\includegraphics[width=5cm,height=0.3cm]{zxjiadanbaojunshu_dis.png}}}\vskip6pt}  &\hspace*{-3.573946226cm}\raisebox{-0.45ex}{\includegraphics{triangle.pdf}} & \begin{minipage}{4.60cm}\begin{center}{偏低 }\end{center} \end{minipage} \\
\hline
\parbox[c]{\hsize}{\vskip6pt 韦荣氏球菌属\\分解乳酸为乙酸和丙酸，调节肠道pH，极少数菌种可能引起炎症 \vskip6pt} & \parbox[c]{\hsize}{\vskip6pt\centerline{\raisebox{-1.5ex}{\includegraphics[width=5cm,height=0.3cm]{zxxiaoweirongqiujunshu_dis.png}}}\vskip6pt}  &\hspace*{-3.56343743cm}\raisebox{-0.45ex}{\includegraphics{triangle.pdf}} & \begin{minipage}{4.60cm}\begin{center}{偏低\\ \bahao 不利于产生有益物质及调节肠道pH值 }\end{center} \end{minipage} \\
\hline
\parbox[c]{\hsize}{\vskip6pt 乳酸杆菌属\\肠道益生菌，能够生成乳酸，抑制有害菌及炎症，调节肠道环境 \vskip6pt} & \parbox[c]{\hsize}{\vskip6pt\centerline{\raisebox{-1.5ex}{\includegraphics[width=5cm,height=0.3cm]{zxruganjunshu_dis.png}}}\vskip6pt}  &\hspace*{-3.81751205cm}\raisebox{-0.45ex}{\includegraphics{triangle.pdf}} & \begin{minipage}{4.60cm}\begin{center}{偏低\\ \bahao 不利于抑制肠道炎症及调节肠道环境 }\end{center} \end{minipage} \\
\hline
\parbox[c]{\hsize}{\vskip6pt 肠球菌属\\多数为共生菌，产生乳酸。少数菌种可能导致多种感染症状 \vskip6pt} & \parbox[c]{\hsize}{\vskip6pt\centerline{\raisebox{-1.5ex}{\includegraphics[width=5cm,height=0.3cm]{zxchangqiujunshu_dis.png}}}\vskip6pt}  &\hspace*{-3.502131638cm}\raisebox{-0.45ex}{\includegraphics{triangle.pdf}} & \begin{minipage}{4.60cm}\begin{center}{偏低\\ \bahao 不利于肠道菌群平衡 }\end{center} \end{minipage} \\
\hline
\parbox[c]{\hsize}{\vskip6pt 嗜胆菌属\\共生菌，可能与长期高脂高蛋白饮食有关 \vskip6pt} & \parbox[c]{\hsize}{\vskip6pt\centerline{\raisebox{-1.5ex}{\includegraphics[width=5cm,height=0.3cm]{zxshidanjunshu_dis.png}}}\vskip6pt}  &\hspace*{-3.463162142cm}\raisebox{-0.45ex}{\includegraphics{triangle.pdf}} & \begin{minipage}{4.60cm}\begin{center}{偏低 }\end{center} \end{minipage} \\
\hline
\parbox[c]{\hsize}{\vskip6pt 沙雷氏菌属\\多数为致病菌，导致尿路感染、腹膜炎等多种医院内感染 \vskip6pt} & \parbox[c]{\hsize}{\vskip6pt\centerline{\raisebox{-1.5ex}{\includegraphics[width=5cm,height=0.3cm]{zxshaleishijunshu_dis.png}}}\vskip6pt}  &\hspace*{-3.600634544cm}\raisebox{-0.45ex}{\includegraphics{triangle.pdf}} & \begin{minipage}{4.60cm}\begin{center}{偏低 }\end{center} \end{minipage} \\
\hline
\parbox[c]{\hsize}{\vskip6pt 放线菌属\\多数为共生菌，少数在免疫力弱时可引起内源性感染、龋齿等 \vskip6pt} & \parbox[c]{\hsize}{\vskip6pt\centerline{\raisebox{-1.5ex}{\includegraphics[width=5cm,height=0.3cm]{zxfangxianjunshu_dis.png}}}\vskip6pt}  &\hspace*{-3.470946878cm}\raisebox{-0.45ex}{\includegraphics{triangle.pdf}} & \begin{minipage}{4.60cm}\begin{center}{偏低\\ \bahao 不利于肠道菌群平衡 }\end{center} \end{minipage} \\
\hline
\parbox[c]{\hsize}{\vskip6pt 毛杆菌属\\肠道共生菌，发酵葡萄糖产生乳酸及少量乙酸和丁酸 \vskip6pt} & \parbox[c]{\hsize}{\vskip6pt\centerline{\raisebox{-1.5ex}{\includegraphics[width=5cm,height=0.3cm]{zxlachnobacterium_dis.png}}}\vskip6pt}  &\hspace*{-2.33468219cm}\raisebox{-0.45ex}{\includegraphics{triangle.pdf}} & \begin{minipage}{4.60cm}\begin{center}{偏低\\ \bahao 不利于肠道菌群平衡 }\end{center} \end{minipage} \\
\hline
\parbox[c]{\hsize}{\vskip6pt 葡萄球菌属\\多数为共生菌，分解葡萄糖等产酸。少数可引起感染、食物中毒等 \vskip6pt} & \parbox[c]{\hsize}{\vskip6pt\centerline{\raisebox{-1.5ex}{\includegraphics[width=5cm,height=0.3cm]{zxputaoqiujunshu_dis.png}}}\vskip6pt}  &\hspace*{-3.582041714cm}\raisebox{-0.45ex}{\includegraphics{triangle.pdf}} & \begin{minipage}{4.60cm}\begin{center}{偏低\\ \bahao 不利于肠道菌群平衡 }\end{center} \end{minipage} \\
\hline
\parbox[c]{\hsize}{\vskip6pt 芽孢杆菌属\\多数为共生菌，降解复杂碳水化合物。少数可引起感染、食物中毒等 \vskip6pt} & \parbox[c]{\hsize}{\vskip6pt\centerline{\raisebox{-1.5ex}{\includegraphics[width=5cm,height=0.3cm]{zxyabaoganjunshu_dis.png}}}\vskip6pt}  &\hspace*{-3.53444636cm}\raisebox{-0.45ex}{\includegraphics{triangle.pdf}} & \begin{minipage}{4.60cm}\begin{center}{偏低\\ \bahao 不利于肠道菌群平衡 }\end{center} \end{minipage} \\
\hline
\parbox[c]{\hsize}{\vskip6pt 阿德勒氏菌属\\帮助分解食物中的异黄酮产生雌马酚，有利于肠道及人体健康 \vskip6pt} & \parbox[c]{\hsize}{\vskip6pt\centerline{\raisebox{-1.5ex}{\includegraphics[width=5cm,height=0.3cm]{zxadlercreutzia_dis.png}}}\vskip6pt}  &\hspace*{-4.121226812cm}\raisebox{-0.45ex}{\includegraphics{triangle.pdf}} & \begin{minipage}{4.60cm}\begin{center}{偏低\\ \bahao 不利于肠道及人体健康 }\end{center} \end{minipage} \\
\hline
\parbox[c]{\hsize}{\vskip6pt 嗜蛋白胨菌属\\利用蛋白胨作为主要能源，产生丁酸，与多组织及血液感染有关 \vskip6pt} & \parbox[c]{\hsize}{\vskip6pt\centerline{\raisebox{-1.5ex}{\includegraphics[width=5cm,height=0.3cm]{zxshidongjunshu_dis.png}}}\vskip6pt}  &\hspace*{-3.441406514cm}\raisebox{-0.45ex}{\includegraphics{triangle.pdf}} & \begin{minipage}{4.60cm}\begin{center}{偏低\\ \bahao 不利于产生有益物质 }\end{center} \end{minipage} \\
\hline
\parbox[c]{\hsize}{\vskip6pt 克雷伯氏菌属\\多为致病菌，可能导致肺炎、尿路感染、软组织感染、菌血症等 \vskip6pt} & \parbox[c]{\hsize}{\vskip6pt\centerline{\raisebox{-1.5ex}{\includegraphics[width=5cm,height=0.3cm]{zxkeleibojunshu_dis.png}}}\vskip6pt}  &\hspace*{-3.772000328cm}\raisebox{-0.45ex}{\includegraphics{triangle.pdf}} & \begin{minipage}{4.60cm}\begin{center}{偏低 }\end{center} \end{minipage} \\
\hline
\parbox[c]{\hsize}{\vskip6pt 脱硫弧菌属\\产生硫化氢，刺激肠道产生炎症反应，不利于肠道健康 \vskip6pt} & \parbox[c]{\hsize}{\vskip6pt\centerline{\raisebox{-1.5ex}{\includegraphics[width=5cm,height=0.3cm]{zxtuoliuhujunshu_dis.png}}}\vskip6pt}  &\hspace*{-2.192120228cm}\raisebox{-0.45ex}{\includegraphics{triangle.pdf}} & \begin{minipage}{4.60cm}\begin{center}{偏低 }\end{center} \end{minipage} \\
\hline
\parbox[c]{\hsize}{\vskip6pt 爱格士氏菌属\\多为致病菌，与溃疡性结肠炎、肛门脓肿、菌血症等疾病有关 \vskip6pt} & \parbox[c]{\hsize}{\vskip6pt\centerline{\raisebox{-1.5ex}{\includegraphics[width=5cm,height=0.3cm]{zxaigetejunshu_dis.png}}}\vskip6pt}  &\hspace*{-4.111154762cm}\raisebox{-0.45ex}{\includegraphics{triangle.pdf}} & \begin{minipage}{4.60cm}\begin{center}{偏低 }\end{center} \end{minipage} \\
\hline
\parbox[c]{\hsize}{\vskip6pt 乳球菌属\\发酵葡萄糖产生乳酸，调节肠道pH，抑制有害菌的生长 \vskip6pt} & \parbox[c]{\hsize}{\vskip6pt\centerline{\raisebox{-1.5ex}{\includegraphics[width=5cm,height=0.3cm]{zxruqiujunshu_dis.png}}}\vskip6pt}  &\hspace*{-3.547123946cm}\raisebox{-0.45ex}{\includegraphics{triangle.pdf}} & \begin{minipage}{4.60cm}\begin{center}{偏低\\ \bahao 不利于调节肠道pH及抑制有害菌生长 }\end{center} \end{minipage} \\
\hline
\parbox[c]{\hsize}{\vskip6pt 真杆菌属\\产生短链脂肪酸等有益物质，少数可能与炎症、菌血症等相关 \vskip6pt} & \parbox[c]{\hsize}{\vskip6pt\centerline{\raisebox{-1.5ex}{\includegraphics[width=5cm,height=0.3cm]{zxzhenganjunshu_dis.png}}}\vskip6pt}  &\hspace*{-3.499966832cm}\raisebox{-0.45ex}{\includegraphics{triangle.pdf}} & \begin{minipage}{4.60cm}\begin{center}{偏低\\ \bahao 不利于产生有益物质 }\end{center} \end{minipage} \\
\hline
\parbox[c]{\hsize}{\vskip6pt 埃希氏菌属\\多数为共生菌，产生维生素K，少数与腹泻、尿道感染、败血症等相关 \vskip6pt} & \parbox[c]{\hsize}{\vskip6pt\centerline{\raisebox{-1.5ex}{\includegraphics[width=5cm,height=0.3cm]{zxaixishijunshu_dis.png}}}\vskip6pt}  &\hspace*{-4.27342607cm}\raisebox{-0.45ex}{\includegraphics{triangle.pdf}} & \begin{minipage}{4.60cm}\begin{center}{偏低\\ \bahao 不利于肠道菌群平衡 }\end{center} \end{minipage} \\
\hline
\parbox[c]{\hsize}{\vskip6pt 别样杆菌属\\共生菌，降解多糖产生琥珀酸和乙酸，可能与长期高脂饮食相关 \vskip6pt} & \parbox[c]{\hsize}{\vskip6pt\centerline{\raisebox{-1.5ex}{\includegraphics[width=5cm,height=0.3cm]{zxalistipes_dis.png}}}\vskip6pt}  &\hspace*{-3.820688792cm}\raisebox{-0.45ex}{\includegraphics{triangle.pdf}} & \begin{minipage}{4.60cm}\begin{center}{偏低\\ \bahao 不利于产生有益物质 }\end{center} \end{minipage} \\
\hline
\parbox[c]{\hsize}{\vskip6pt 丁酸弧菌属\\降解多种碳水化合物，产生丁酸，参与蛋白质和脂类代谢 \vskip6pt} & \parbox[c]{\hsize}{\vskip6pt\centerline{\raisebox{-1.5ex}{\includegraphics[width=5cm,height=0.3cm]{zxdingsuanhujunshu_dis.png}}}\vskip6pt}  &\hspace*{-3.415527446cm}\raisebox{-0.45ex}{\includegraphics{triangle.pdf}} & \begin{minipage}{4.60cm}\begin{center}{偏低\\ \bahao 不利于产生有益物质及物质代谢 }\end{center} \end{minipage} \\
\hline
\parbox[c]{\hsize}{\vskip6pt 弯曲杆菌属\\多数菌种为致病菌，可引起弯曲菌病，表现为严重腹泻或痢疾综合征 \vskip6pt} & \parbox[c]{\hsize}{\vskip6pt\centerline{\raisebox{-1.5ex}{\includegraphics[width=5cm,height=0.3cm]{zxwanqujunshu_dis.png}}}\vskip6pt}  &\hspace*{-3.69573611cm}\raisebox{-0.45ex}{\includegraphics{triangle.pdf}} & \begin{minipage}{4.60cm}\begin{center}{偏低 }\end{center} \end{minipage} \\
\hline
\end{longtable}

\noindent
\tikz\draw[green2,fill=green2](0,0) circle[radius=4pt]; 区域表示正常人群的一般水平 \tikz\draw[darkblue,fill=darkblue](0,0) circle[radius=4pt]; 区域表示仍在正常水平，但较一般水平偏高或偏低 %\tikz\draw[red,fill=red](0,0) circle[radius=4pt]; 而红色表示为有风险
}

\bigskip
{\centering {\includegraphics{jieshi.pdf}}}

\end{document}
